\section*{Bab \ref{ch:b3:inorganic-materials} --- Material Anorganik dan Nanomaterial}

\begin{solution}{exo:b3:inorganic-materials:1}
$x^2 \propto t$: tiga kali lebih tebal memerlukan waktu sembilan kali lebih lama, \qty{90}{h}.
\end{solution}

\begin{solution}{exo:b3:inorganic-materials:2}
Hidrolisis: $\ce{Si(OC2H5)4 + 4H2O -> Si(OH)4 + 4C2H5OH}$.\\
Kondensasi: $\ce{2Si(OH)4 -> Si2O(OH)6 + H2O}$.\\
Keseluruhan: $\ce{Si(OC2H5)4 + 2H2O -> SiO2 + 4C2H5OH}$.
\end{solution}

\begin{solution}{exo:b3:inorganic-materials:3}
Pembentuk: $\ce{SiO2}$, $\ce{B2O3}$. Pengubah: $\ce{Na2O}$, $\ce{CaO}$. Setiap ion oksida dari $\ce{Na2O}$ memutus satu jembatan Si--O--Si
menjadi dua ujung Si--O$^-$, yang masing-masing berpasangan dengan sebuah $\ce{Na+}$: jaringannya terpotong, lelehannya
mengalir lebih mudah dan meleleh pada suhu yang lebih rendah.
\end{solution}

\begin{solution}{exo:b3:inorganic-materials:4}
Dua belas segi lima (seperti pada setiap \omterm{def:b3:inorganic-materials:carbon}{fulerena}). $V = 70$, $E = 3V/2 = 105$ ikatan, $F = 2 - V + E = 37$
muka, sehingga ada $37 - 12 = 25$ segi enam.
\end{solution}

\begin{solution}{exo:b3:inorganic-materials:5}
$\sqrt2\,(r_B + r_O) = 1.4142 \times \qty{200.5}{pm} = \qty{283.5}{pm}$. $\ce{SrTiO3}$: $284/283.5 = 1.00$, perovskit kubus ideal.
$\ce{BaTiO3}$: $301/283.5 = 1.06$, titanium terlalu kecil untuk oktahedronnya dan bergeser dari pusat: \omterm{def:b3:inorganic-materials:ferroelectric}{feroelektrik}.
$\ce{CaTiO3}$: $274/283.5 = 0.97$, kalsium terlalu kecil untuk rongganya: oktahedron-oktahedronnya miring, simetrinya
lebih rendah.
\end{solution}

\begin{solution}{exo:b3:inorganic-materials:6}
Si/Al $= 106/86 = 1.23$. $M = 86 \times 23.0 + 86 \times 27.0 + 106 \times 28.1 + 384 \times 16.0 = \qty{13422.6}{g/mol}$; 43 $\ce{Ca^2+}$ per
rumus: $43/\qty{13422.6}{g/mol} = \qty{3.20}{mmol/g}$.
\end{solution}

\begin{solution}{exo:b3:inorganic-materials:7}
Jarak tetangga terdekat $d = a/\sqrt2 = \qty{288.4}{pm}$. Gugusan dengan $n$ kulit berdiameter $(2n + 1)d$:
$2n + 1 \approx 5/0.2884 = 17.3$, sehingga $n = 8$ (\qty{4.9}{nm}). $N(8) = 2057$ atom, $10 \times 64 + 2 = 642$ di permukaan:
31\,\%.
\end{solution}

\begin{solution}{exo:b3:inorganic-materials:8}
$h^2/(8\mu R^2)$ dengan $\mu = \qty{9.109e-32}{kg}$: \qty{0.94}{eV} pada \qty{2.0}{nm} dan \qty{0.42}{eV} pada \qty{3.0}{nm}. Emisi pada
$1.74 + 0.94 = \qty{2.68}{eV}$, \qty{463}{nm} (biru), dan $1.74 + 0.42 = \qty{2.16}{eV}$, \qty{574}{nm} (kuning): \omterm{def:b3:inorganic-materials:quantum-dot}{titik kuantum} yang lebih kecil
memancarkan cahaya yang lebih biru.
\end{solution}

\begin{solution}{exo:b3:inorganic-materials:9}
$(10, 10)$: kursi, $n - m = 0$, logam. $(10, 0)$: zigzag, 10 bukan kelipatan 3, \omterm{def:b3:solid-state:classes}{semikonduktor}.
$(9, 0)$: zigzag, logam. $(12, 8)$: bukan keduanya (kiral), $n - m = 4$, \omterm{def:b3:solid-state:classes}{semikonduktor}.
\end{solution}

\begin{solution}{exo:b3:inorganic-materials:10}
$M = 4 \times 65.4 + 13 \times 16.0 + 24 \times 12.0 + 12 \times 1.0 = \qty{769.6}{g/mol}$; $a^3 = (\qty{25.82e-8}{cm})^3 = \qty{1.721e-20}{cm^3}$;
$\rho = 8 \times 769.6/(\num{6.022e23} \times \num{1.721e-20}) = \qty{0.594}{g.cm^{-3}}$. Satu lapangan seluas \qty{7140}{m^2}: satu gram memuat
$3000/7140 = 0.42$ lapangan, dan satu sendok teh berisi beberapa gram melebihi satu lapangan penuh.
\end{solution}

\begin{solution}{exo:b3:inorganic-materials:11}
$\dfrac{10n^2 + 2}{(10n^3 + 15n^2 + 11n + 3)/3} = \dfrac{30n^2 + 6}{10n^3 + 15n^2 + 11n + 3} \to \dfrac{30n^2}{10n^3} = \dfrac3n$. Untuk $n = 8$: nilai eksaknya
$642/2057 = 0.31$, dibandingkan dengan $3/8 = 0.375$; pendekatan itu melebih-lebihkan gugusan
kecil, yang kulit-kulit dalamnya tidak dapat diabaikan.
\end{solution}

\begin{solution}{exo:b3:inorganic-materials:12}
$2E = 3V$, $2E = 5p + 6h + 7s$, $F = p + h + s$. Euler: $V - E + F = 2$, dengan $V = 2E/3$: $F - E/3 = 2$, sehingga
$p + h + s - (5p + 6h + 7s)/6 = 2$, yaitu $(p - s)/6 = 2$: $p - s = 12$. Setiap segi tujuh memerlukan satu segi lima
tambahan.
\end{solution}

\begin{solution}{pb:b3:inorganic-materials:1}
\textbf{1.} $\mathrm{Na_{12}Al_{12}Si_{12}O_{48}}$: $12 \times 23.0 + 12 \times 27.0 + 12 \times 28.1 + 48 \times 16.0 = \qty{1705.2}{g/mol}$.

\textbf{2.} $1705.2 + 27 \times 18.0 = \qty{2191.2}{g/mol}$.

\textbf{3.} Si/Al $= 12/12 = 1$.

\textbf{4.} Tidak ada tautan Al--O--Al; dengan Si/Al $= 1$, setiap Si memiliki empat tetangga Al dan setiap
Al empat tetangga Si: berselang-seling secara ketat.

\textbf{5.} $-12$, satu per aluminium.

\textbf{6.} $486.0/2191.2 = 22.2$\,\%.

\textbf{7.} $\ce{2Na+}(\text{zeolit}) + \ce{Ca^2+}(\text{aq}) \rightleftharpoons \ce{Ca^2+}(\text{zeolit}) + \ce{2Na+}(\text{aq})$.

\textbf{8.} Enam (dua belas muatan).

\textbf{9.} $6/\qty{1705.2}{g/mol} = \qty{3.52}{mmol/g}$.

\textbf{10.} $6/\qty{2191.2}{g/mol} = \qty{2.74}{mmol/g}$.

\textbf{11.} $\qty{3.52}{mmol/g} \times \qty{100.1}{g/mol} = \qty{352}{mg}$ $\ce{CaCO3}$ per gram.

\textbf{12.} $\qty{2.0}{mmol/L} \times \qty{15}{L} = \qty{30}{mmol}$.

\textbf{13.} $\qty{30}{mmol}/\qty{2.74}{mmol/g} = \qty{11}{g}$.

\textbf{14.} $\qty{10.9}{g}/0.60 = \qty{18}{g}$.

\textbf{15.} $\qty{60}{mmol}$ $\ce{Na+}$, $\qty{0.060}{mol} \times \qty{23.0}{g/mol} = \qty{1.4}{g}$.

\textbf{16.} Fosfat dalam air limbah menyuburkan alga di sungai dan danau; ledakan alga
dan pembusukannya menghabiskan oksigen terlarut (eutrofikasi).

\textbf{17.} $\ce{Mg^2+}$ yang lebih kecil memegang molekul-molekul airnya lebih erat; ion itu harus melepaskannya
untuk melewati jendela dan mencapai situs-situs, dan hal itu lambat pada suhu pencucian.

\textbf{18.} Ion tanpa selubung berdiameter sekitar \qty{0.20}{nm}, setengah lebar jendela.

\textbf{19.} Ion terhidrat lebih besar daripada jendela: ion itu harus melepaskan sebagian
selubung airnya untuk lewat, dan oksigen-oksigen kerangka menggantikan tempat
air itu.

\textbf{20.} Kerangka itu bermuatan negatif dan hanya mempertukarkan kation; anion
besar dengan rantai panjang ditolak dan tidak dapat melewati jendela, sehingga tetap
berada di dalam air, tempat anion itu bekerja.

\textbf{21.} Dalam keadaan terdehidrasi, \omterm{def:b3:inorganic-materials:zeolite}{zeolit} menyerap air dengan kuat ke dalam sangkar-sangkarnya; hanya
molekul yang lebih kecil daripada jendela yang dapat masuk, sehingga \omterm{def:b3:inorganic-materials:zeolite}{zeolit} memisahkan molekul menurut ukurannya:
sebuah ayakan dalam skala molekul.

\textbf{22.} Ion-ion $\ce{K+}$ yang lebih besar menempati jendela-jendela dan mempersempitnya: hanya molekul
yang lebih kecil, seperti air, yang dapat diterima.

\textbf{23.} Isomer \textit{para} yang ramping dan linear muat di dalam saluran dan berdifusi cepat; isomer
\textit{orto} dan \textit{meta} yang lebih lebar berdifusi lambat, tetap tinggal, dan berisomerisasi.

\textbf{24.} Secara teoretis, \omterm{def:b3:inorganic-materials:zeolite}{zeolit} A anhidrat dapat mengambil \qty{3.52}{mmol} $\ce{Ca^2+}$ per gram.
\end{solution}
